\documentclass[preview]{standalone}

\usepackage{amsmath}
\usepackage{amssymb}
\usepackage{stellar}
\usepackage{definitions}

\begin{document}

\id{psicologia-emozioni-fisiologia}
\genpage

\section{Emozioni e fisiologia}

\begin{snippetdefinition}{emozione-definition}{Emozione}
    Un'\emph{emozione} è un'esperienza complessa costituita da cambiamenti
    corporei e mentali, che includono attivazione fisiologica, sentimenti,
    processi cognitivi, espressioni visibili e reazioni comportamentali
    specifiche, attivate in risposta a situazioni percepite come personalmente
    significative.
\end{snippetdefinition}

\begin{snippetexample}{felicita-example}{Felicità}
    Quando si prova felicità, l'attivazione fisiologica può determinare un
    leggero aumento del battito cardiaco. Il comportamento manifesto può essere
    espressivo, come nel caso di un sorriso, oppure orientato all'azione, come
    quando si abbraccia una persona cara.
\end{snippetexample}

\begin{snippet}{funzioni-emozioni}
    Le emozioni influenzano l'azione e preparano l'organismo a rispondere a
    situazioni rilevanti. La paura davanti a un pericolo improvviso può
    paralizzare oppure indurre alla fuga; altre emozioni orientano il legame
    sociale, la comunicazione e la valutazione degli eventi.
\end{snippet}

\begin{snippet}{attivazione-fisiologica-emozioni}
    Durante una forte emozione, il cuore batte più rapidamente, la respirazione
    aumenta e i muscoli si irrigidiscono. Questi cambiamenti sono coordinati da
    molteplici processi fisiologici che agiscono sull'intero organismo.
\end{snippet}

\begin{snippetdefinition}{sistema-nervoso-autonomo-emozioni-definition}{Sistema nervoso autonomo nelle emozioni}
    Il \emph{sistema nervoso autonomo}, o \emph{SNA}, prepara il corpo alla
    risposta emotiva attraverso l'azione coordinata del sistema simpatico e del
    sistema parasimpatico.
\end{snippetdefinition}

\begin{snippet}{simpatico-parasimpatico-emozioni}
    Con una stimolazione moderata e piacevole si attiva soprattutto il sistema
    parasimpatico. Con una stimolazione più intensa, come paura o rabbia, il
    sistema simpatico prepara il corpo alla reazione di emergenza, favorendo il
    rilascio di adrenalina e noradrenalina e aumentando pressione sanguigna,
    sudorazione e disponibilità di energia.
\end{snippet}

\begin{snippetdefinition}{amigdala-definition}{Amigdala}
    L'\emph{amigdala} è una struttura cerebrale coinvolta nella decodifica delle
    informazioni provenienti dagli organi sensoriali e nell'attribuzione di
    significato a esperienze emotivamente salienti, soprattutto negative.
\end{snippetdefinition}

\begin{snippet}{cervello-emozioni}
    La ricerca sulle emozioni si è concentrata su ipotalamo, sistema limbico e
    amigdala. Le emozioni positive e negative non producono semplicemente
    risposte opposte nella stessa area corticale, ma coinvolgono reti cerebrali
    distribuite e possono mostrare differenze individuali.
\end{snippet}

\end{document}
